\section{Port Knocking}

Technik um versteckte Dienste zu finden.
Ports in bestimmter Reihenfolge ansprechen:

\begin{lstlisting}[language=bash]
# Einzelne Ports scannen
nmap -p PORT target_ip

# Port Knocking mit hping3
hping3 -S target_ip -p PORT1
hping3 -S target_ip -p PORT2
hping3 -S target_ip -p PORT3
\end{lstlisting}

TCP-Flags in der Antwort:
\begin{description}
	\item[RA (Reset+ACK):] Port geschlossen
	\item[SA (SYN+ACK):] Port offen
\end{description}

\section{Forensisches Image erstellen}

Ablauf: leeres Image erstellen → partitionieren → formatieren → Daten kopieren → EWF erzeugen:

\begin{lstlisting}[language=bash]
# Leeres Image erstellen
dd if=/dev/zero of=usb.img bs=1M count=512

# Partitionieren
fdisk usb.img
# -> n (new), p (primary), Typ aendern: t, 07 (NTFS)

# Loop Device und formatieren
losetup --partscan --find --show usb.img
mkfs.ntfs /dev/loop0p1

# Mounten und Daten kopieren
mount /dev/loop0p1 /mnt
rsync -av source/ /mnt/

# EWF-Image erstellen
ewfacquire /dev/loop0
# Case-Nummer, Examiner, Beschreibung angeben
\end{lstlisting}

\section{Netzwerkadapter-Analyse}

Über Windows Registry:

\begin{lstlisting}[language=bash]
regripper -p nic2 -r /tmp/SYSTEM
\end{lstlisting}

Zeigt: Adapter-Name, DHCP-Server, IP-Adresse, Subnetzmaske, Gateway, DNS.
Forensisch relevant: verbundene Netzwerke (z.B. Fritz.box → Heimnetzwerk).

\section{FUSE-Treiber}

Filesystem in Userspace — ermöglicht Zugriff auf verschiedene Dateisysteme:
\begin{itemize}
	\item \texttt{ntfs-3g} — NTFS lesen/schreiben
	\item \texttt{xmount} — FUSE-basiert für Image-Konvertierung
\end{itemize}
